AI scientist systems that autonomously analyze data and write research papers are scaling rapidly, yet no prior work has evaluated whether they correctly recover a known data generating process.
We introduce the \livesalmon---a controlled benchmark inspired by the classic dead salmon fMRI study \citep{bennett2009neural}---in which three frontier \llms (\gptmodel, \claudemodel, \geminimodel) each analyze a synthetic dataset with a known formula ($Y = 2X_1 + 0.5X_2 - 1.5X_3 + 0.8 \cdot X_4 X_5 + \epsilon$) and a null-signal control dataset, repeated three times per model (18 papers total).
We find that \aiscientists are competent but methodologically narrow: all 9 live-salmon papers identify the strong linear predictors $X_1$ and $X_3$ (100\%), and 67\% detect the weaker $X_2$.
However, \textbf{zero out of 9 papers detect the $X_4 \times X_5$ interaction term}, which accounts for 8\% of total variance---a systematic blind spot arising from uniform reliance on correlation analysis and linear regression.
On the dead salmon data, all 9 papers correctly report null findings with zero false positives, contrasting sharply with the spurious activations in the original fMRI study.
These results reveal that the collective ``AI scientific literature'' would mischaracterize the true model as purely linear, highlighting a critical gap between statistical conservatism and analytical creativity in current \aiscientist systems.
